%! Characteristics of Polynomial Functions II — Unit 3, Lesson 3.5 — Guided Notes & Practice (Answer Key)
% MIRROR of ../notes/main.tex (Units 3+ packet-first, LESSON_SHAPE.md §0a). Identical except:
%   \usepackage{algebra2-boxes} -> \usepackage{algebra2-key}; the page-header title;
%   \blank{W} -> \ans{...} (no wider than its blank); \writelines{2} -> two \ansline{}\\ lines.
% Every \begin{work} block and every figure is BYTE-IDENTICAL. \ans{} is TEXT mode.
% NO teachernote here — teacher prose goes in the lesson plan.
\documentclass[12pt]{article}
\usepackage{algebra2-article}
\usepackage{algebra2-key}   % pulls in -boxes; never load both
\newcommand{\rung}[1]{{\color{goldacc}\ifcase#1\or$\star$\or$\star\star$\or$\star\star\star$\fi}}
% ── Pre-drawn polynomial graphs (never a sketch task). One grid style everywhere: unit squares,
% axes through the origin, every zero and every turning point dotted and labelled.
% \pgrid{xmin}{xmax}{ymin}{ymax}[x-name][y-name]-style axes; \pcurve clips to the window.
\newcommand{\pgrid}[6]{%
  \draw[step=1, linegray!55, very thin] (#1,#3) grid (#2,#4);
  \draw[->, semithick, charcoal] (#1,0)--({#2+0.5},0) node[right, font=\tiny]{$#5$};
  \draw[->, semithick, charcoal] (0,#3)--(0,{#4+0.5}) node[above, font=\tiny]{$#6$};}
\newcommand{\pcurve}[6]{%
  \begin{scope}\clip (#1,#3) rectangle (#2,#4);
    \draw[very thick, forest, domain=#5, samples=140, smooth] plot (\x,{#6});
  \end{scope}}
% \zr[label position]{x}{label} — a zero on the x-axis;  \tp[label position]{x}{y}{label} — a
% turning point (or other key point).
\newcommand{\zr}[3][below]{\fill[charcoal] (#2,0) circle (2.4pt)
  node[#1, font=\tiny, inner sep=1.2pt, fill=white, fill opacity=0.85, text opacity=1]{#3};}
\newcommand{\tp}[4][above]{\fill[navy] (#2,#3) circle (2.4pt)
  node[#1, font=\tiny, inner sep=1.2pt, fill=white, fill opacity=0.85, text opacity=1]{#4};}
% The lesson's graphs, by name (the deck draws the same ones). Optional argument = scale.
% A(x) = (x+2)(x-1)^2 — Round 1's display (zeros annotated) and Round 2's Show me (turning points).
\newcommand{\gAdef}[1][0.42]{\begin{tikzpicture}[scale=#1]
  \pgrid{-3}{3}{-2}{5}{x}{y}\pcurve{-3}{3}{-2}{5}{-2.6:2.3}{(\x+2)*(\x-1)*(\x-1)}
  \zr[below right]{-2}{$-2$: crosses}\zr[above right]{1}{$1$: touches}\end{tikzpicture}}
\newcommand{\gA}[1][0.42]{\begin{tikzpicture}[scale=#1]
  \pgrid{-3}{3}{-2}{5}{x}{y}\pcurve{-3}{3}{-2}{5}{-2.6:2.3}{(\x+2)*(\x-1)*(\x-1)}
  \zr[below left]{-2}{$-2$}\tp[right]{-1}{4}{$(-1,4)$}\tp[below right]{1}{0}{$(1,0)$}\end{tikzpicture}}
% p(x) = -1/4 (x+3)(x+1)^2(x-2) — Round 1 rung 3 (zeros only).
\newcommand{\gP}[1][0.42]{\begin{tikzpicture}[scale=#1]
  \pgrid{-4}{3}{-2}{5}{x}{y}\pcurve{-4}{3}{-2}{5}{-3.6:2.6}{-0.25*(\x+3)*(\x+1)*(\x+1)*(\x-2)}
  \zr[below left]{-3}{$-3$}\zr[below]{-1}{$-1$}\zr[below right]{2}{$2$}\end{tikzpicture}}
% W(x) = 1/4 (x^2-4)^2 — Round 2 example.
\newcommand{\gW}[1][0.42]{\begin{tikzpicture}[scale=#1]
  \pgrid{-3}{3}{-1}{5}{x}{y}\pcurve{-3}{3}{-1}{5}{-3:3}{0.25*(\x*\x-4)*(\x*\x-4)}
  \tp[below]{-2}{0}{$(-2,0)$}\tp[above right]{0}{4}{$(0,4)$}\tp[below]{2}{0}{$(2,0)$}\end{tikzpicture}}
% C(x) = (x+1)(x-2)^2 — Round 2 rung 2.
\newcommand{\gC}[1][0.42]{\begin{tikzpicture}[scale=#1]
  \pgrid{-2}{4}{-2}{5}{x}{y}\pcurve{-2}{4}{-2}{5}{-1.6:3.4}{(\x+1)*(\x-2)*(\x-2)}
  \zr[below left]{-1}{$-1$}\tp[above right]{0}{4}{$(0,4)$}\tp[below]{2}{0}{$(2,0)$}\end{tikzpicture}}
% M(x) = 4 - 1/4 (x^2-4)^2 — Round 2 rung 3.
\newcommand{\gM}[1][0.42]{\begin{tikzpicture}[scale=#1]
  \pgrid{-4}{4}{-2}{5}{x}{y}\pcurve{-4}{4}{-2}{5}{-3.2:3.2}{4-0.25*(\x*\x-4)*(\x*\x-4)}
  \tp[above left]{-2}{4}{$(-2,4)$}\tp[above right]{2}{4}{$(2,4)$}\tp[below]{0}{0}{$(0,0)$}
  \zr[below left]{-2.828}{$-2\sqrt2$}\zr[below right]{2.828}{$2\sqrt2$}\end{tikzpicture}}
% E(x) = (x+1)^2(x-2) — Round 3 example.
\newcommand{\gE}[1][0.42]{\begin{tikzpicture}[scale=#1]
  \pgrid{-3}{3}{-5}{2}{x}{y}\pcurve{-3}{3}{-5}{2}{-2.4:2.7}{(\x+1)*(\x+1)*(\x-2)}
  \tp[above left]{-1}{0}{$(-1,0)$}\tp[below right]{1}{-4}{$(1,-4)$}\zr[above left]{2}{$2$}\end{tikzpicture}}
% K(x) = -(x+2)^2(x-1) — Round 3 rung 1.
\newcommand{\gK}[1][0.42]{\begin{tikzpicture}[scale=#1]
  \pgrid{-4}{2}{-2}{5}{x}{y}\pcurve{-4}{2}{-2}{5}{-3.6:1.6}{-1*(\x+2)*(\x+2)*(\x-1)}
  \tp[below]{-2}{0}{$(-2,0)$}\tp[above right]{0}{4}{$(0,4)$}\zr[below left]{1}{$1$}\end{tikzpicture}}
% S(x) = x^3(x-4)/9 — Round 3 rung 2 (one turning point; flat crossing at 0).
\newcommand{\gS}[1][0.42]{\begin{tikzpicture}[scale=#1]
  \pgrid{-2}{5}{-4}{3}{x}{y}\pcurve{-2}{5}{-4}{3}{-1.9:4.6}{(\x*\x*\x*(\x-4))/9}
  \zr[above left]{0}{$0$}\zr[above left]{4}{$4$}\tp[below]{3}{-3}{$(3,-3)$}\end{tikzpicture}}
% T(h) = (h-2)^2(h-5)/4, 0 <= h <= 6 — Round 3 rung 3 (the riding barn).
\newcommand{\gT}[1][0.42]{\begin{tikzpicture}[scale=#1]
  \pgrid{0}{6}{-5}{4}{h}{T}\pcurve{0}{6}{-5}{4}{0:6}{(\x-2)*(\x-2)*(\x-5)/4}
  \tp[right]{0}{-5}{$(0,-5)$}\tp[above]{2}{0}{$(2,0)$}\tp[below]{4}{-1}{$(4,-1)$}
  \zr[below right]{5}{$5$}\tp[left]{6}{4}{$(6,4)$}\end{tikzpicture}}
% R(x) — Find the mistake (turning points only).
\newcommand{\gR}[1][0.42]{\begin{tikzpicture}[scale=#1]
  \pgrid{-5}{3}{-3}{5}{x}{y}\pcurve{-5}{3}{-3}{5}{-5.2:3.2}{(3*\x*\x*\x+9*\x*\x-27*\x-17)/16}
  \tp[above]{-3}{4}{$(-3,4)$}\tp[below]{1}{-2}{$(1,-2)$}\end{tikzpicture}}
% H(x) = -(x-1)^2(x+2) — Closing Practice.
\newcommand{\gH}[1][0.42]{\begin{tikzpicture}[scale=#1]
  \pgrid{-3}{3}{-5}{2}{x}{y}\pcurve{-3}{3}{-5}{2}{-2.6:2.4}{-1*(\x-1)*(\x-1)*(\x+2)}
  \zr[above left]{-2}{$-2$}\tp[below left]{-1}{-4}{$(-1,-4)$}\tp[above right]{1}{0}{$(1,0)$}\end{tikzpicture}}

\begin{document}

\pageheader{Unit 3, Lesson 3.5}{Guided Notes \& Practice --- Answer Key}
% NAMESTRIP: no \namedateperiod here — the student writes their name once, on the cover.

\vspace{0.10in}

% ═══ ROUND 1 — ZEROS AND MULTIPLICITY ═══════════════════════════════════════
\boxguard[20]
\begin{notesbox}{Round 1 --- Zeros and multiplicity}
\small
\begin{minipage}[t]{0.60\linewidth}
\vspace{0pt}
\textbf{Zero.} A \textbf{zero} of $f$ is an input $x$ with $f(x)=0$. On the graph it is an
$x$-intercept; in factored form, each factor $(x-r)$ gives the zero $x=r$.

\vspace{4pt}
\textbf{Multiplicity.} The \textbf{multiplicity} of a zero is the number of times its factor
appears --- the exponent on that factor. At a zero of \textbf{odd} multiplicity the graph
\textbf{crosses} the $x$-axis; at a zero of \textbf{even} multiplicity it \textbf{touches} the
axis and turns back (it bounces).
\end{minipage}\hfill
\begin{minipage}[t]{0.36\linewidth}
\vspace{0pt}\centering
\gAdef[0.40]\\[1pt]
{\scriptsize $A(x)=(x+2)^{1}(x-1)^{2}$}
\end{minipage}

\vspace{6pt}
\textbf{Example.} List each zero of $f(x)=(x+3)^2(x-1)(x-4)^3$, its multiplicity, and whether
the graph crosses or touches there.
\begin{work}
(x+3)^2 &:\ x=-3,\ \text{multiplicity } 2\ (\text{even}) \Rightarrow \text{touches} \\
(x-1)   &:\ x=1,\ \text{multiplicity } 1\ (\text{odd}) \Rightarrow \text{crosses} \\
(x-4)^3 &:\ x=4,\ \text{multiplicity } 3\ (\text{odd}) \Rightarrow \text{crosses} \\
\text{degree} &= 2+1+3=6 \quad \text{(the multiplicities add to the degree)}
\end{work}

\vspace{4pt}
\textbf{Your turn.} \textit{Start at \rung{1} and climb.} List each zero, its multiplicity, and
crosses or touches.

\rung{1}\; $f(x)=(x-5)(x+1)^2$
\begin{work}
(x-5)   &:\ x=5,\ \text{multiplicity } 1\ (\text{odd}) \Rightarrow \text{crosses} \\
(x+1)^2 &:\ x=-1,\ \text{multiplicity } 2\ (\text{even}) \Rightarrow \text{touches}
\end{work}

\rung{2}\; $g(x)=-x^2(x-3)^3(x+6)$
\begin{work}
x^2     &:\ x=0,\ \text{multiplicity } 2\ (\text{even}) \Rightarrow \text{touches} \\
(x-3)^3 &:\ x=3,\ \text{multiplicity } 3\ (\text{odd}) \Rightarrow \text{crosses} \\
(x+6)   &:\ x=-6,\ \text{multiplicity } 1\ (\text{odd}) \Rightarrow \text{crosses} \\
-1      &:\ \text{not a factor with a zero --- it only flips the graph}
\end{work}

\begin{minipage}[t]{0.60\linewidth}
\vspace{0pt}
\rung{3}\; The graph of a polynomial $p$ is shown. Read its zeros.

\vspace{6pt}
Zeros: \ans{$x=-3,\ -1,\ 2$}

\vspace{6pt}
Which zero has \emph{even} multiplicity? $x=$ \ans{$-1$}

\vspace{6pt}
How do you know? \ans{it touches, not crosses}

\vspace{6pt}
Least possible degree of $p$: \ans{$4$ \ ($1+2+1$)}
\end{minipage}\hfill
\begin{minipage}[t]{0.36\linewidth}
\vspace{0pt}\centering
\gP[0.40]
\end{minipage}

\vspace{8pt}
\textbf{Show me} (options on the screen). Where does $f(x)=(x+2)^3(x-4)^2$ cross, and where does
it touch? \quad My answer: \ans{1}
\end{notesbox}

% ═══ ROUND 2 — TURNING POINTS AND EXTREMA ═══════════════════════════════════
\boxguard[24]
\begin{notesbox}{Round 2 --- Turning points and extrema}
\small
\textbf{Turning point.} A \textbf{turning point} is a point where the graph changes from
increasing to decreasing, or from decreasing to increasing. A polynomial of degree $n$ has
\textbf{at most $n-1$} turning points.

\vspace{4pt}
\textbf{Relative vs.\ absolute.} The $y$-value at a turning point is a \textbf{relative
maximum} if it is the highest point \emph{near there}, and a \textbf{relative minimum} if it is
the lowest point near there. An \textbf{absolute maximum} (or \textbf{absolute minimum}) is the
greatest (or least) $y$-value of the \emph{whole} function. An odd-degree polynomial has no
absolute maximum or minimum. An even-degree polynomial has one absolute extreme value: an
absolute minimum when $a>0$, an absolute maximum when $a<0$.

\vspace{6pt}
\begin{minipage}[t]{0.62\linewidth}
\vspace{0pt}
\textbf{Example.} The graph of $f(x)=\tfrac14(x^2-4)^2$, a degree-4 polynomial, is shown. Count
its turning points and name every relative and absolute extremum.
\begin{work}
\text{degree } 4 &\Rightarrow \text{at most } 3 \text{ turning points; it has } 3 \\
\text{relative max} &= 4 \text{ at } x=0 \\
\text{relative min} &= 0 \text{ at } x=-2 \text{ and } x=2 \\
\text{absolute min} &= 0 \quad (\text{even degree, } a>0) \\
\text{absolute max} &: \text{none --- the arms rise past } 4
\end{work}
\end{minipage}\hfill
\begin{minipage}[t]{0.34\linewidth}
\vspace{0pt}\centering
\gW[0.40]
\end{minipage}

\vspace{6pt}
\textbf{Your turn.} \textit{Start at \rung{1} and climb.}

\rung{1}\; A polynomial has degree $5$. At most how many turning points? \ans{4}

\vspace{4pt}
\hspace*{1.3em}Another polynomial's graph has $3$ turning points. Its least possible degree is
\ans{4}.

\vspace{6pt}
\begin{minipage}[t]{0.62\linewidth}
\vspace{0pt}
\rung{2}\; The graph of $C(x)=(x+1)(x-2)^2$ is shown.

\vspace{6pt}
relative max \ans{4} at $x=$ \ans{0}

\vspace{6pt}
relative min \ans{0} at $x=$ \ans{2}

\vspace{6pt}
absolute max and min: \ans{none --- odd degree}
\end{minipage}\hfill
\begin{minipage}[t]{0.34\linewidth}
\vspace{0pt}\centering
\gC[0.40]
\end{minipage}

\vspace{6pt}
\begin{minipage}[t]{0.62\linewidth}
\vspace{0pt}
\rung{3}\; The graph of a degree-4 polynomial $M$ is shown.

\vspace{6pt}
absolute max \ans{4} at $x=$ \ans{$-2$, $2$}

\vspace{6pt}
relative min \ans{0} at $x=$ \ans{0}

\vspace{6pt}
Is $0$ the absolute minimum? \ans{no --- the arms fall forever}
\end{minipage}\hfill
\begin{minipage}[t]{0.34\linewidth}
\vspace{0pt}\centering
\gM[0.36]
\end{minipage}

\vspace{8pt}
\textbf{Show me} (options on the screen). On the Round 1 graph of $A(x)$, what is the turning
point $(-1,4)$? \quad My answer: \ans{2}
\end{notesbox}

% ═══ ROUND 3 — INTERVALS, DOMAIN AND RANGE (the crux) ════════════════════════
\boxguard[24]
\begin{notesbox}{Round 3 --- Intervals, domain, and range}
\small
\textbf{Increasing and decreasing.} A function is \textbf{increasing} on an interval if its
$y$-values rise as $x$ moves left to right, and \textbf{decreasing} if they fall. We name the
interval by its \textbf{$x$-values}, in interval notation, with parentheses at the turning points.

\vspace{4pt}
\textbf{Positive and negative.} $f$ is \textbf{positive} where its graph is above the $x$-axis
($f(x)>0$) and \textbf{negative} where it is below ($f(x)<0$); these intervals start and stop at
zeros.

\vspace{4pt}
\textbf{Domain and range.} The \textbf{domain} of every polynomial is $(-\infty,\infty)$. The
\textbf{range} of an odd-degree polynomial is $(-\infty,\infty)$; an even-degree polynomial's
range stops at its absolute extremum, such as $[-3,\infty)$.

\vspace{6pt}
\begin{minipage}[t]{0.62\linewidth}
\vspace{0pt}
\textbf{Example.} Use the graph of $E(x)=(x+1)^2(x-2)$.
\begin{work}
\text{increasing} &:\ (-\infty,-1)\cup(1,\infty) \\
\text{decreasing} &:\ (-1,1) \quad (x\text{-values only}) \\
\text{positive}   &:\ (2,\infty) \\
\text{negative}   &:\ (-\infty,-1)\cup(-1,2) \quad (E(-1)=0) \\
\text{domain}     &:\ (-\infty,\infty),\quad \text{range}:\ (-\infty,\infty)
\end{work}
\end{minipage}\hfill
\begin{minipage}[t]{0.34\linewidth}
\vspace{0pt}\centering
\gE[0.40]
\end{minipage}

\vspace{6pt}
\textbf{Your turn.} \textit{Start at \rung{1} and climb.} Answer in interval notation.

\vspace{4pt}
\begin{minipage}[t]{0.62\linewidth}
\vspace{0pt}
\rung{1}\; The graph of $K(x)=-(x+2)^2(x-1)$ is shown.

\vspace{6pt}
increasing on \ans{$(-2,\,0)$}

\vspace{6pt}
domain \ans{$(-\infty,\,\infty)$}
\end{minipage}\hfill
\begin{minipage}[t]{0.34\linewidth}
\vspace{0pt}\centering
\gK[0.36]
\end{minipage}

\vspace{6pt}
\begin{minipage}[t]{0.62\linewidth}
\vspace{0pt}
\rung{2}\; The graph of $S(x)=\tfrac19x^3(x-4)$ is shown.

\vspace{6pt}
decreasing on \ans{$(-\infty,\,3)$}

\vspace{6pt}
negative on \ans{$(0,\,4)$}

\vspace{6pt}
range \ans{$[-3,\,\infty)$}
\end{minipage}\hfill
\begin{minipage}[t]{0.34\linewidth}
\vspace{0pt}\centering
\gS[0.36]
\end{minipage}

\vspace{6pt}
\begin{minipage}[t]{0.62\linewidth}
\vspace{0pt}
\rung{3}\; On a cold night the temperature in the riding barn is $T$ degrees Celsius, $h$ hours
after midnight, for $0\le h\le6$.

\vspace{6pt}
When was the barn warming up? \ans{$(0,\,2)\cup(4,\,6)$}

\vspace{6pt}
When was it below freezing ($T<0$)? \ans{$[0,\,2)\cup(2,\,5)$}
\end{minipage}\hfill
\begin{minipage}[t]{0.34\linewidth}
\vspace{0pt}\centering
\gT[0.36]
\end{minipage}

\vspace{8pt}
\begin{minipage}[t]{0.62\linewidth}
\vspace{0pt}
\textbf{Find the mistake.} Student A read the graph of $R$ and wrote:

\vspace{3pt}
\hspace*{1em}``$R$ is increasing on $(-\infty,4)\cup(-2,\infty)$ and decreasing on $(-2,4)$.''

\vspace{3pt}
Which part is wrong? Explain, then fix it.
\end{minipage}\hfill
\begin{minipage}[t]{0.34\linewidth}
\vspace{0pt}\centering
\gR[0.32]
\end{minipage}
\par
\ansline{Student A used the turning points' $y$-values ($4$ and $-2$); intervals use $x$-values.}\\
\ansline{Increasing on $(-\infty,-3)\cup(1,\infty)$; decreasing on $(-3,1)$.}\\
\end{notesbox}

% ═══ CLOSING PRACTICE ═══════════════════════════════════════════════════════
\boxguard[24]
\begin{notesbox}{Closing Practice --- score yourself out of 4}
\small
\begin{minipage}[t]{0.64\linewidth}
\vspace{0pt}
\begin{enumerate}[leftmargin=1.5em, itemsep=8pt]
  \item (Round 1) The zeros of $f(x)=(x-6)^2(x+5)^3$, with multiplicity and crosses/touches:

        \ans{$6$: mult.\ 2, touches; \ $-5$: mult.\ 3, crosses}
  \item (Round 2) A degree-6 polynomial has at most \ans{5} turning points. Does it have
        an absolute extremum? \ans{yes}
  \item (Round 3) Using the graph of $H$: $H$ is increasing on \ans{$(-1,\,1)$}.
  \item (Mixed) For $H$: the relative minimum is \ans{$-4$}, and $H$ is positive on
        \ans{$(-\infty,\,-2)$}.
\end{enumerate}
\end{minipage}\hfill
\begin{minipage}[t]{0.32\linewidth}
\vspace{0pt}\centering
\gH[0.40]\\[1pt]
{\scriptsize $H(x)=-(x-1)^2(x+2)$}
\end{minipage}

\vspace{8pt}
\textbf{Bonus (+1).} During evening study hall the number of students in the library, $L(t)$,
follows a degree-3 model, $t$ hours after 7 p.m. It has a relative maximum of $40$ students at
$t=1$ and a relative minimum of $12$ students at $t=2.5$. On what interval is $L$ decreasing?
\ans{$(1,\,2.5)$}
\end{notesbox}

\end{document}
